% !TeX program = xelatex
\documentclass[UTF8,11pt,a4paper,oneside,openany,linespread=1]{ctexbook}
\usepackage[margin=25mm]{geometry}
\usepackage{amsmath,amssymb,bm,booktabs,array,enumitem}
\usepackage{setspace}
\setstretch{1.28}
\setlength{\parskip}{3pt}
\setlength{\jot}{5pt}
\raggedbottom
\usepackage{xcolor}
\definecolor{Ink}{HTML}{19364A}
\definecolor{Teal}{HTML}{227C81}
\definecolor{Rust}{HTML}{9B483A}
\usepackage[most]{tcolorbox}
\tcbset{enhanced,breakable,boxrule=.45pt,arc=1mm,
  left=3mm,right=3mm,top=2mm,bottom=2mm,
  before skip=9pt,after skip=9pt,fonttitle=\bfseries}
\newtcolorbox{conceptbox}[1]{title={概念：#1},colframe=Ink,colback=Ink!3}
\newtcolorbox{methodbox}[1]{title={方法卡：#1},colframe=Teal,colback=Teal!3}
\newtcolorbox{pitfallbox}{title=我踩过的坑,colframe=Rust,colback=Rust!3}
\newtcbtheorem[number within=chapter]{notetheorem}{定理}
  {colframe=Ink,colback=Ink!3}{thm}
\newcommand{\R}{\mathbb{R}}
\newcommand{\vect}[1]{\bm{#1}}
\newcommand{\dd}{\mathop{}\!\mathrm{d}}
\begin{document}
\setcounter{chapter}{3}
\chapter{自定义环境与常用命令}

\section{列表该不该编号}
并列的材料用 itemize，有先后顺序的步骤用 enumerate。
\begin{itemize}
  \item 定理及其适用条件；
  \item 一个完整例题；
  \item 一条值得回查的计算提醒。
\end{itemize}
\begin{enumerate}[label=\arabic*.,leftmargin=*]
  \item 检查条件和定义域。
  \item 选择计算方法。
  \item 代回原式检查结果。
\end{enumerate}
我不手敲编号。中间插入一步以后，自动编号比挨个改数字省事。

\section{把知识框定义一次}
\begin{conceptbox}{原函数}
在一个区间上，若 \(F'=f\)，就称 \(F\) 是 \(f\) 的原函数。
\end{conceptbox}
\begin{methodbox}{先检查，再计算}
先写适用条件，再给出计算过程，最后代回检查。
\end{methodbox}
\begin{pitfallbox}
两个原函数之差为常数，需要在同一个区间上理解。
\end{pitfallbox}
\begin{notetheorem}{原函数差为常数}{constant}
若 \(F\) 与 \(G\) 是区间上同一函数的原函数，则 \(F-G\) 为常数。
\end{notetheorem}

\verb|\newtcolorbox| 后第一组花括号是环境名称，\verb|[1]| 表示定义的环境接收一个必填参数。
样式里的 \verb|#1| 是使用时填写的标题。固定标题的踩坑框不接收参数。
我把共同设置放进 tcbset，再给不同环境指定标题和颜色。
\begin{center}
\begin{tabular}{ll}
\toprule
参数 & 控制内容 \\
\midrule
colframe / colback & 边框颜色与正文底色 \\
colbacktitle / coltitle & 标题条底色与标题文字颜色 \\
boxrule / arc & 边框粗细与圆角 \\
left / right / top / bottom & 框内留白 \\
before skip / after skip & 框前后的空隙 \\
breakable & 是否允许框跨页 \\
\bottomrule
\end{tabular}
\end{center}
长推导用普通正文更易读。我用 \verb|\newtcbtheorem| 为定理自动编号。

\section{用命令统一记号}
我把常用数集、向量和微分记号定义在导言区：
\[
  D\subseteq\R,\qquad\vect{x}\in\R^n,\qquad
  \int_0^1x^2\,\dd x=\frac13.
\]
没有参数的命令代替固定记号，有参数的命令封装一类操作。
以后改变向量风格，只改 vect 的定义，不用全书替换字母。

newcommand 遇到已有名字会报错，renewcommand 则要求原定义存在。
先查重名来自哪里，再决定换名或重定义，别为了让日志安静而盲目覆盖。
\end{document}
